% TeamBench v2. Venue-neutral preamble on purpose: no style file names a venue.
\documentclass[11pt]{article}
\usepackage[margin=1in]{geometry}
\usepackage{booktabs}
\usepackage{amsmath,amssymb}
\usepackage{graphicx}
\usepackage{xcolor}
\usepackage[hidelinks]{hyperref}
\usepackage{natbib}

\newcommand{\tb}{\textsc{TeamBench}}
\newcommand{\core}{\textsc{TeamBench-Core}}
\newcommand{\todo}[1]{\textcolor{red}{[#1]}}
\newcommand{\cn}{\todo{CITATION NEEDED}}

\title{What a Multi-Agent Coding Benchmark Measures When It Measures Nothing}
\author{Anonymous}
\date{}

\begin{document}
\maketitle

\begin{abstract}
We show that a multi-agent coding benchmark can report a complete set of
results while being unable to distinguish any two submissions, and we give a
protocol that detects this before any model is run. Applying it to our own
prior benchmark, 31 of 76 graders could not separate a correct submission from
an empty one, because the test that defines each bug is the one the upstream
pull request adds and the task ships the repository state from before it;
another 40 tasks shipped that test inside the workspace, handing the agent the
acceptance criterion. Two grader checks could not be satisfied by any
submission at all, including the maintainers' own merged fix. Across the
corpus, 303 of 374 checks pass on an untouched workspace, so 81\% of every
reported partial score was credit for not damaging the repository. We rebuild
the corpus by checking each task out from upstream at its recorded base commit
and injecting the held-out tests at grade time, admit only tasks that fail
without the reference fix and pass with it, and verify the result against seven
admission gates. On the 48 tasks that survive, we then re-run the comparison the
original paper could not settle: at matched compute a Planner/Executor/Verifier
team solved one task in 144 paired comparisons that a single agent failed, the
deficit widens as compute grows ($-0.104 \to -0.188$), and team runs damage the
workspace three times as often. On the scale the original paper used, none of
this is visible: a sevenfold compute increase moves the reported score by 1.5
points.
\end{abstract}

\section{Introduction}

A benchmark reports numbers whether or not it can measure anything. A grader
that awards the same score to every submission still produces a leaderboard, a
pass rate, and an ablation table. Nothing in the usual reporting pipeline
distinguishes a task that is hard from a task that is inert.

This paper is a report on finding that condition in our own work, the protocol
that now detects it, and what the central experiment looks like once the corpus
is repaired. We take the position that a benchmark paper owes the reader
evidence that its tasks can discriminate, on the same footing as evidence about
the systems it evaluates.

\paragraph{The failure.} The GitHub-derived half of \tb{} was built by copying,
for each merged bug-fix pull request, the files that pull request touched, at
the commit before it. Two consequences follow, and they hide each other.

First, the copied workspace is a fragment of a repository. Its test suite cannot
be collected: \texttt{conftest.py} is absent on 541 of the 553 tasks that ship a
test file, packages are present only in the pieces the diff touched, and
test-only dependencies are missing. The grader's central check fails for reasons
unrelated to any submission.

Second, and underneath it, the test that \emph{defines} the bug is the test the
pull request adds, and the commit before it does not have that test. Of the 229
tasks whose reference patch adds a test, 189 ship the test file without the
added test, so the grader cannot separate any two submissions; the other 40 ship
the added test inside the workspace, so the agent can read the acceptance
criterion. The only tasks that could discriminate were the ones leaking the
answer.

\paragraph{Contributions.}
\begin{itemize}
\item \textbf{A validity protocol for agentic benchmarks.} Stage the task,
      grade an empty submission, grade the reference fix, and require that the
      first fails and the second passes. We run it on every task and report the
      count that survives, not the count that exists.
\item \textbf{Evidence that graders can be inert, quantified.} 31 of 76 graders
      were blind to the bug; two check families could not be satisfied by any
      submission, including the reference fix; 81\% of the checks are guards
      that pass on an untouched workspace.
\item \textbf{\core{}, 48 tasks that provably discriminate,} rebuilt by
      checking out upstream at the recorded base commit and holding the tests
      out until grade time, and validated against seven admission gates.
\item \textbf{The compute-matched comparison, re-run.} At equal turn budget a
      three-role team is not better than a single agent, the gap widens with
      compute, and the team damages the workspace more often. We also show that
      the scale used in the prior work hides the effect entirely.
\end{itemize}

\section{How a grader becomes inert}
\label{sec:inert}

\todo{FIGURE 1: the three workspace states. (a) upstream at base\_sha, no added
test: the grader cannot see the bug. (b) the vendored fragment for 40 tasks,
added test present: the agent can read the criterion. (c) ours: base\_sha
workspace, test injected at grade time.}

\paragraph{Setting.} Each task derives from a merged pull request that fixes a
bug. The task's workspace is the repository at the parent commit; the reference
solution is the merged diff. A grader runs a checklist and reports a pass flag
and a partial score.

\paragraph{The held-out test.} A bug-fix pull request usually adds the test that
demonstrates the bug. At the parent commit that test does not exist. A workspace
built from the parent commit therefore contains a test suite that passes, and
will keep passing whatever the agent does, unless the agent happens to reproduce
the maintainers' test. Measured on our corpus: applying the maintainers' own
merged fix changed the reported score on 3 of 153 tasks.

The repair is the one SWE-bench uses \cn: hold the tests out of the workspace
and inject them when grading. We precompute the post-patch content of each test
file at stage time and copy it over the workspace before the grader runs, rather
than applying a patch. Copying is deterministic and overwrites whatever the
agent wrote there, so weakening a test cannot raise a score.

\paragraph{Checks no submission can satisfy.} Two check families failed with the
reference fix in place.

The import check loaded a package-internal module as a detached top-level
module:
\begin{quote}\small
\texttt{spec = spec\_from\_file\_location('mod', src\_path); exec\_module(...)}
\end{quote}
Every relative import inside such a file raises
\texttt{attempted relative import with no known parent package}. It failed with
the reference fix applied on 37 of 59 staged tasks. Because the grader's overall
pass requires every check, those tasks were unsolvable by construction.

The second is subtler. Graders invoked \texttt{pytest} on a single file and
inherited the repository's own \texttt{addopts}. For a project with a coverage
gate this is fatal: with the fix applied and the held-out tests injected,
\texttt{arrow} reports \texttt{223 passed} and then
\texttt{FAIL Required test coverage of 99\% not reached}. Every test passes and
the grader records a failure.

\paragraph{Guards in the score.} 303 of the 374 checks across the admitted
corpus pass on an untouched workspace. They are not defective; they are guards
that detect deleted tests, wiped files and hardcoded shortcuts. The defect is
that they sat in the score. An empty submission earned a mean partial score of
0.815. We move guards out of the score rather than removing them: a submission
that trips one is reported inadmissible, and the score is computed over the
checks that fail on that task's pristine workspace. Discriminative status is a
property of a (task, check) pair and never of a check's name: ``source modules
import without error'' is a guard on one task and the entire point of another.

\section{\core{}}
\label{sec:core}

\paragraph{Construction.} For each task we check the upstream repository out at
its recorded base commit, install it in an isolated environment, and ask two
questions, following SWE-bench's fail-to-pass criterion \cn:
\begin{align*}
G_{\text{fail}} &: \text{the tests the patch adds fail without the source fix} \\
G_{\text{pass}} &: \text{the same tests pass with it}
\end{align*}
Both must hold. We record the resolved dependency set so staging never
re-resolves against whatever the package index serves later.

\paragraph{The ceiling is the honest denominator.} Of 631 tasks carrying a
repository and base commit, 229 (36\%) have a reference patch that adds a test.
The remaining 402 cannot be verified by a test signal, because the upstream pull
request never added one. Reporting the admitted count against 631 rather than
against 229 would make a sound corpus look broken.

\paragraph{Validation.} Each candidate is graded three ways: pristine, pristine
with the held-out tests injected, and with the full reference fix. A task is
admitted iff the second fails and the third passes.

\begin{table}[t]
\centering\small
\begin{tabular}{lrr}
\toprule
 & usable & graded \\
\midrule
before the import-check repair          & 10 & 59 \\
after it                                & 33 & 66 \\
after the coverage-gate repair          & \textbf{48} & \textbf{76} \\
\bottomrule
\end{tabular}
\caption{Tasks that provably discriminate, over the same candidate set. The
first two numbers were set by defects in the grader, not by the corpus. In the
final round, 31 of the 76 graders could not see the bug at all until the
held-out tests were injected.}
\label{tab:validation}
\end{table}

\core{} is the 48 admitted tasks, drawn from 12 upstream repositories. Six of
seven admission gates pass on essentially every one of them: an empty submission
fails all 48, each has a non-empty discriminative set, and the upstream
reference scores 1.0 on 47 of 48.

The seventh gate fails, and we report it as a failure. It asks that at least
half of a grader's checks be discriminative; the measured median is 0.14.
Satisfying it would mean deleting anti-cheat checks, so we do not adjust it.
What it now measures is granularity rather than validity, and the number is
itself a result: with a median of one discriminative check per task, \core{} is
effectively pass/fail. The graded partial credit reported in the prior version
of this work was 81\% guard mass.

\section{Compute-matched coordination}
\label{sec:budget}

\paragraph{The confound.} Every comparison in the prior version ran the
conditions at different compute: a single agent at 20 model turns, a three-role
team at up to 140. Any difference so measured is confounded with a sevenfold
budget gap. We give every condition one shared turn allowance, consumed by
whichever role is acting.

\paragraph{Setup.} 48 tasks $\times$ 2 conditions $\times$ 3 budgets, one model,
seed 0. Solo receives the full specification; Full Team is
Planner/Executor/Verifier with enforced information partition.

\begin{table}[t]
\centering\small
\begin{tabular}{lrrrrr}
\toprule
budget & condition & pass & raw mean & disc.\ mean & guard violations \\
\midrule
20  & Solo      & 4/48 \ (8\%)  & 0.826 & 0.104 & 1 \\
20  & Full Team & 0/48 \ (0\%)  & 0.803 & 0.000 & 1 \\
60  & Solo      & 5/48 \ (10\%) & 0.839 & 0.149 & 2 \\
60  & Full Team & 0/48 \ (0\%)  & 0.793 & 0.031 & 4 \\
140 & Solo      & 10/48 (21\%)  & 0.841 & 0.267 & 3 \\
140 & Full Team & 2/48 \ (4\%)  & 0.779 & 0.080 & 8 \\
\bottomrule
\end{tabular}
\caption{288 runs at matched compute. ``raw'' is the score as the grader emits
it, guards included; ``disc.'' is over the checks that fail on an untouched
workspace. A guard violation is a submission that deleted tests or wiped a file.}
\label{tab:budget}
\end{table}

\paragraph{At equal compute the team does not win.} Paired within task and
budget, the team solved one task the single agent failed, across 144
comparisons. Of 19 discordant pairs, 18 favour the single agent (exact
two-sided sign test, $p = 7.6\times10^{-5}$).

\paragraph{More compute makes it relatively worse.} The mean discriminative
deficit is $-0.104$, $-0.118$ and $-0.188$ at 20, 60 and 140 turns. ``The team
needs more compute'' is therefore not available as an explanation.

\paragraph{Teams damage the workspace more.} Guard violations rise
$1 \to 4 \to 8$ of 48 for the team against $1 \to 2 \to 3$ for the single agent.
At 140 turns 17\% of team runs deleted a test or wiped a file, against 6\% of
solo runs. Three roles means three hands editing a repository without seeing
what the others did. \todo{This is a mechanism claim; support it with the
failure taxonomy or soften it.}

\paragraph{The scale decides what is visible.}
\begin{center}\small
\begin{tabular}{llrrr}
\toprule
scale & condition & 20 & 60 & 140 \\
\midrule
raw  & Solo & 0.826 & 0.839 & 0.841 \\
raw  & Team & 0.803 & 0.793 & 0.779 \\
disc.& Solo & 0.104 & 0.149 & 0.267 \\
disc.& Team & 0.000 & 0.031 & 0.080 \\
\bottomrule
\end{tabular}
\end{center}
On the raw scale a sevenfold compute increase moves the single agent by 1.5
points and the conditions sit on top of each other, because 81\% of the checks
are guards and everything is compressed against the ceiling. That is the scale
the prior version reported on, which is why that comparison could not have
settled the question either way.

\section{Related work}
\todo{TO BE WRITTEN. Every citation below must be fetched programmatically
before use; none are written from memory.}
\begin{itemize}
\item SWE-bench and the fail-to-pass criterion. \cn
\item Agentic coding benchmarks and their construction. \cn
\item Multi-agent LLM systems and role decomposition. \cn
\item \textbf{Must cite early:} Gu \& Kim, Nature Machine Intelligence, same
      first author, currently uncited. \cn
\item CooperBench, Silo-Bench, Protocol Validity, The Verifier Tax. \cn
\item Benchmark validity and construct validity in ML evaluation. \cn
\end{itemize}

\section{Limitations}

\paragraph{One model, one seed.} The budget curve uses a single model at seed 0.
The sign test is over task-to-task variation, not over model or seed variation.
\todo{A second model is running; fold in when complete.}

\paragraph{\core{} is small and narrow.} 48 tasks from 12 repositories, all
Python. The ceiling of 229 is itself limited by how often upstream maintainers
add a test with their fix. Repositories that build their own compiled
extensions (165 tasks, ceiling 53) and those whose installation pulls a deep
learning framework (94 tasks, ceiling 23) are excluded and would require
per-task container images.

\paragraph{Effectively pass/fail.} With a median of one discriminative check per
task, \core{} does not support graded partial credit. We report the
discriminative mean, but it is close to a pass rate.

\paragraph{A residual leak.} Task identifiers end in the upstream pull request
number. Renaming would invalidate every recorded result and run path, so we
record it instead. It is weaker than a URL: the model must resolve the reference
unaided. We removed 1061 explicit links to the fix pull request from the
documents the agents read.

\paragraph{Our own protocol found our own defects.} Every count in
Section~\ref{sec:inert} describes a benchmark we built and published. We do not
know how many of these conditions are present in benchmarks we did not build.

\bibliographystyle{plainnat}
\bibliography{references}

\end{document}
